\chapter{Collatz}\label{ch:collatz}

\begin{center}\itshape
For about a month everybody at Yale worked on it, with no result. A similar phenomenon happened when I mentioned it at the University of Chicago. A
joke was made that this problem was part of a conspiracy to slow down mathematical research in the U.S.
\upshape(Shizuo Kakutani, as quoted by Jeffrey Lagarias, 1985)
\end{center}

Take a positive integer. If it is even, halve it; if it is odd, triple it and add one. Repeat. Starting from $6$ you get $6,3,10,5,16,8,4,2,1$, and then
the loop $4,2,1$ for ever. Starting from $7$ the walk takes $16$ steps and climbs to $52$ on the way. Starting from $27$ it takes $111$ steps, climbs to
$9232$, and only then comes down (Figure~\ref{fig:traj27}). The \emph{Collatz conjecture} says that every starting number eventually reaches $1$.

The origin of the problem is, in Lagarias's words, obscure. Lothar Collatz, then a student in Germany, wrote down iteration problems of this kind in his
notebooks in the early 1930s; the one in his notebook dated 1 July 1932 is a related map, a permutation of the integers, about which he asked whether
the cycle containing $8$ is finite. The $3n+1$ problem itself circulated by word of mouth and was known by the early 1950s. Helmut Hasse talked about it
so much that it became ``Hasse's algorithm'', and his visit to Syracuse University gave it the name ``Syracuse problem''; Stanis\l aw Ulam spread it at
Los Alamos, Shizuo Kakutani at Yale and Chicago, and Bryan Thwaites claimed to have found it himself in 1952. Coxeter offered $50$ dollars for a
solution in 1970, Erd\H{o}s $500$, Thwaites $1000$ pounds. Erd\H{o}s is reported to have said that mathematics may not be ready for such problems.

It is the easiest open problem to state in this book and perhaps the hardest to attack. Every number below $2^{71}$ has been checked (David Barina,
2025). What can actually be proved is less, but it is beautiful: if you start from a ``typical'' number, the first steps behave \emph{exactly} like
fair coin tosses, which explains why the walk goes down on average, and gives a theorem of Riho Terras (1976): almost every number eventually drops
below its starting value. A second argument, with continued fractions in disguise, shows that a counterexample of the other kind, a cycle, would have to
be absurdly long.

\begin{think}
Write the walk of $27$ in terms of the steps ``odd'' and ``even''. After an odd step the number is even, so the next step is a halving. On average,
how much does one ``odd step and the halving after it'' multiply the number by? And how many halvings follow on average?
\end{think}

\begin{figure}[ht]
\centering
\fig{ch06-traj27}
\caption{The walk of $27$ under the Collatz map, on a logarithmic scale: $111$ steps, with a maximum of $9232$ at step $77$. The climbs and falls
look like a random walk, and the next sections make that precise.}
\label{fig:traj27}
\end{figure}

\section{The shortcut map and its evidence}

After a step $n\mapsto3n+1$ the number is even, so we merge that step with the halving that follows. The \emph{shortcut map} is
\[
T(n)=\begin{cases}\dfrac n2&n\text{ even},\\[2mm]\dfrac{3n+1}2&n\text{ odd}.\end{cases}
\]
An orbit of the Collatz map $C(n)=\frac n2$ or $3n+1$ reaches $1$ if and only if the corresponding orbit of $T$ does, since the $T$-orbit is the
$C$-orbit with some terms left out. The only cycle known is $1\to2\to1$ for $T$ (that is, $1\to4\to2\to1$ for $C$).

The \emph{stopping time} of $n\ge1$ is $\sigma(n)=\min\{k\ge1:T^k(n)<n\}$, with $\sigma(n)=\infty$ if the orbit never goes below $n$. For instance
$\sigma(1)=\infty$, every even number has $\sigma(n)=1$, and Exercise~\ref{ex:stopsmall} finds the next few cases. The conjecture is a statement about
stopping times.

\begin{proposition}\label{prop:stopping}
Every orbit of $T$ reaches $1$ if and only if $\sigma(n)<\infty$ for every $n\ge2$.
\end{proposition}
\begin{proof}
If the orbit of $n\ge2$ reaches $1<n$, then $\sigma(n)$ is finite. Conversely, assume that every $\sigma(n)$ with $n\ge2$ is finite, and prove by strong
induction on $n$ that the orbit of $n$ reaches $1$. For $n=1$ there is nothing to prove. For $n\ge2$, the number $m=T^{\sigma(n)}(n)$ is a positive
integer smaller than $n$, so its orbit reaches $1$ by induction, and the orbit of $n$ contains that of $m$.
\end{proof}

\paragraph{The record.} The verification of the conjecture has been pushed by many people since the 1960s. David Barina's program, running on graphics cards, verified every number below $2^{68}$ (published in
2021) and then every number below $2^{71}\approx2.4\cdot10^{21}$ (published in 2025). These programs do not follow each orbit to $1$: by
Proposition~\ref{prop:stopping}, it is enough to check that every $n$ drops below itself, and Theorem~\ref{thm:paritybij} below lets such a
program discard whole residue classes modulo $2^k$ at once.

Some walks are long. Among the numbers up to $10^6$, the longest walk under $C$ belongs to $837\,799$, with $524$ steps; the stopping-time records under
$T$ are $27$ ($59$ steps before going below $27$), $703$ ($81$), $10\,087$ ($105$), $35\,655$ ($135$), $270\,271$ ($164$), $362\,343$ ($165$), $381\,727$
($173$) and $626\,331$ ($176$). On average, though, the walks are short and regular: the mean number of $C$-steps to reach $1$, over $n\le10^6$, is
$131.4$. We now explain where this number comes from.

\section{Parity vectors behave like coins}

Write $x_i(n)\in\{0,1\}$ for the parity of $T^i(n)$, and $s_k(n)=x_0(n)+\dots+x_{k-1}(n)$ for the number of odd steps among the first $k$. The
\emph{parity vector} of length $k$ of $n$ is $v_k(n)=(x_0(n),\dots,x_{k-1}(n))\in\{0,1\}^k$. For $27$ it starts $1,1,0,1,1,1,1,1,0,\dots$, because the orbit of $T$ is $27,41,62,31,47,71,107,161,242,\dots$

\begin{lemma}\label{lem:affine}
For $n\ge1$ and $k\ge0$,
\[
T^k(n)=\frac{3^{s_k(n)}n+\rho_k(n)}{2^k},
\]
where $\rho_0=0$ and $\rho_{i+1}=\rho_i$ if $x_i(n)=0$, $\rho_{i+1}=3\rho_i+2^i$ if $x_i(n)=1$. In particular $\rho_k(n)\ge0$ depends only on the parity
vector $v_k(n)$.
\end{lemma}
\begin{proof}
Induction on $k$; the case $k=0$ is clear. Write $s=s_k(n)$ and $\rho=\rho_k(n)$. If $T^k(n)$ is even, $T^{k+1}(n)=\frac{3^sn+\rho}{2^{k+1}}$. If it is odd,
$T^{k+1}(n)=\frac12\bigl(3\cdot\frac{3^sn+\rho}{2^k}+1\bigr)=\frac{3^{s+1}n+3\rho+2^k}{2^{k+1}}$.
\end{proof}

\begin{lemma}\label{lem:shift}
Let $n\ge1$, $k\ge1$ and $m\ge0$. For $0\le i\le k$,
\[
T^i(n+m2^k)=T^i(n)+m\,3^{s_i(n)}\,2^{k-i}.
\]
In particular $n$ and $n+m2^k$ have the same parity vector of length $k$.
\end{lemma}
\begin{proof}
Induction on $i$; the case $i=0$ is clear. Let $i<k$ and suppose the formula holds for $i$. The added term $m3^{s_i(n)}2^{k-i}$ is even, because
$k-i\ge1$, so $T^i(n+m2^k)$ and $T^i(n)$ have the same parity. If both are even, halving gives $T^{i+1}(n)+m3^{s_i(n)}2^{k-i-1}$, and
$s_{i+1}(n)=s_i(n)$. If both are odd, $\frac12\bigl(3T^i(n)+3m3^{s_i(n)}2^{k-i}+1\bigr)=T^{i+1}(n)+m3^{s_i(n)+1}2^{k-i-1}$, and
$s_{i+1}(n)=s_i(n)+1$.
\end{proof}

So the parity vector of length $k$ is a function of $n$ modulo $2^k$. The surprise is that it is a bijection: the first $k$ steps of the walk can be
\emph{anything}, and every pattern happens exactly once in every block of $2^k$ consecutive starting values.

\begin{theorem}[Terras, 1976; Everett, 1977]\label{thm:paritybij}
For every $k\ge1$, the map $n\bmod2^k\mapsto v_k(n)$ is a bijection from $\mathbb Z/2^k\mathbb Z$ onto $\{0,1\}^k$. (For the class of $0$ we use the
representative $2^k$.)
\end{theorem}
\begin{proof}
By Lemma~\ref{lem:shift} the map is well defined. Both sets have $2^k$ elements, so it suffices to show it is injective, by induction on $k$. For $k=1$,
$v_1(n)$ is the parity of $n$. Let $k\ge2$, and let $n<n'$ be positive integers that are not congruent modulo $2^k$. If they are not congruent modulo
$2^{k-1}$, their parity vectors of length $k-1$ already differ, by induction. Otherwise $n'\equiv n+2^{k-1}\pmod{2^k}$, and by Lemma~\ref{lem:shift} we may
replace $n'$ by $n+2^{k-1}$. Lemma~\ref{lem:shift} with $k-1$ in place of $k$ and $m=1$ gives, for $i=k-1$,
\[
T^{k-1}(n+2^{k-1})=T^{k-1}(n)+3^{s_{k-1}(n)}.
\]
The added term is odd, so the parities of $T^{k-1}(n)$ and $T^{k-1}(n')$ differ: the vectors differ in their last entry.
\end{proof}

\begin{corollary}[the exact random model]\label{cor:coins}
Among any $2^k$ consecutive positive integers, each parity vector of length $k$ occurs exactly once. Hence exactly $\binom ks$ of them have $s_k(n)=s$, and
\[
\frac1{2^k}\sum_{n=N+1}^{N+2^k}\log\frac{3^{s_k(n)}}{2^k}=k\Bigl(\frac12\log3-\log2\Bigr)=-0.1438\ldots\,k.
\]
\end{corollary}
\begin{proof}
Consecutive integers run once through the residues modulo $2^k$, so the first claim is Theorem~\ref{thm:paritybij}. The number of vectors with $s$ ones is
$\binom ks$, and $\sum_s\binom ks s=k2^{k-1}$, so the mean of $s_k(n)\log3-k\log2$ is $\frac k2\log3-k\log2$.
\end{proof}

In words: if you choose $n$ at random from a long interval, the parities of its first $k$ iterates are $k$ independent fair coins. This is not a
heuristic. By Lemma~\ref{lem:affine}, $T^k(n)$ is about $\frac{3^{s_k(n)}}{2^k}n$ when $n$ is large compared with $\rho_k(n)$, so $\log T^k(n)$ behaves
like a random walk with steps $\log\frac32$ and $\log\frac12$, each with probability $\frac12$, and the mean step $\log\frac{\sqrt3}2=-0.1438$ is
negative. The walk drifts down.

\paragraph{The heuristic.} If the coins stayed fair until the end of the walk, the orbit of $n$ would need about $\frac{\log n}{0.1438}=6.95\log n$
steps of $T$ to reach $1$. Each step of $T$ costs one step of $C$ if it is a halving and two if it is an odd step, $1.5$ on average, so the prediction is
$10.43\log n$ steps of $C$. Averaged over $n\le10^6$, this predicts $133.7$; the true average is $131.4$. The same model says how rare long climbs are: the
proportion of starting values whose walk stays above its start for $k$ steps decays exponentially in $k$ (Figure~\ref{fig:terras}), which is why the
record holders above are so sparse. It does \emph{not} prove anything about a single $n$: the coins
are fair only for the first $k$ steps of $n$ modulo $2^k$, and a particular $n$ is a single, fixed sequence of coins.

\begin{figure}[ht]
\centering
\fig{ch06-tree}
\caption{The backward tree of $1$ to depth $13$: the $61$ numbers whose Collatz walk reaches $1$ in at most $13$ steps, each at the height of its number
of steps. Every $m$ has the predecessor $2m$ (grey edges), and $m\equiv4\pmod6$ also has the odd predecessor $\frac{m-1}3$ (red). The level sizes are
$1,1,1,1,1,2,2,4,4,6,6,8,10,14$.}
\label{fig:tree}
\end{figure}

\section{Terras's theorem}

The coin model gives a real theorem as soon as we can count the parity vectors that go down. We need the classical estimate for the tail of the binomial
distribution, in the form of a pure counting statement.

\begin{lemma}[Chernoff bound]\label{lem:chernoff}
Let $\frac12<a<1$ and $c(a)=\log2+a\log a+(1-a)\log(1-a)$. Then $c(a)>0$, and for every $k\ge1$
\[
\#\bigl\{v\in\{0,1\}^k:\ v\text{ has at least }ak\text{ ones}\bigr\}=\sum_{s\ge ak}\binom ks\le2^ke^{-c(a)k}.
\]
\end{lemma}
\begin{proof}
For every $t>0$, the terms with $s\ge ak$ satisfy $1\le e^{t(s-ak)}$, so
\[
\sum_{s\ge ak}\binom ks\le e^{-tak}\sum_{s=0}^k\binom kse^{ts}=e^{-tak}(1+e^t)^k.
\]
Choose $e^t=\frac a{1-a}>1$. Then $1+e^t=\frac1{1-a}$ and $e^{-ta}=\bigl(\frac{1-a}a\bigr)^a$, so the bound is
$\bigl(a^{-a}(1-a)^{-(1-a)}\bigr)^k=2^ke^{-c(a)k}$. Finally $c(\frac12)=0$ and $c'(a)=\log\frac a{1-a}>0$ for $a>\frac12$, so $c(a)>0$.
\end{proof}

For the Collatz problem the relevant threshold is $a=\frac{\log2}{\log3}=0.6309$: the walk is below its start after $k$ steps, roughly, when
$3^{s_k}<2^k$, that is when $s_k<ak$. Then $c(a)=0.03468\ldots$ Exercise~\ref{ex:chernoffsharp} shows that this exponent cannot be improved.

\begin{theorem}[Terras, 1976]\label{thm:terras}
The set of $n$ with $\sigma(n)<\infty$ has natural density $1$: almost every positive integer eventually goes below itself.
\end{theorem}
\begin{proof}
Fix $k\ge1$ and let $V_k$ be the set of vectors $v\in\{0,1\}^k$ such that $3^{s_j(v)}<2^j$ for some $j\le k$, where $s_j(v)$ is the number of ones among
the first $j$ entries of $v$. For $v\in V_k$ let $j(v)$ be the least such $j$.

Let $n$ be a positive integer with $v_k(n)=v\in V_k$, and $j=j(v)$, $s=s_j(v)$. By Lemma~\ref{lem:affine}, $T^j(n)=\frac{3^sn+\rho_j(n)}{2^j}$, and this
is less than $n$ if and only if $(2^j-3^s)n>\rho_j(n)$. Here $2^j-3^s\ge1$ and $\rho_j(n)$ depends only on $v$. So, with
\[
B_k=\max_{v\in V_k}\frac{\rho_{j(v)}(v)}{2^{j(v)}-3^{s_{j(v)}(v)}},
\]
every $n>B_k$ with $v_k(n)\in V_k$ has $\sigma(n)\le k$. By Theorem~\ref{thm:paritybij}, the $n$ with $v_k(n)\in V_k$ form $|V_k|$ residue classes
modulo $2^k$, so for $x\ge1$
\[
\#\{n\le x:\sigma(n)<\infty\}\ge|V_k|\Bigl(\frac x{2^k}-1\Bigr)-B_k,
\]
hence
\[
\liminf_{x\to\infty}\frac1x\#\{n\le x:\sigma(n)<\infty\}\ge\frac{|V_k|}{2^k}.
\]
A vector not in $V_k$ has in particular $3^{s_k(v)}\ge2^k$, that is $s_k(v)\ge ak$ with $a=\frac{\log2}{\log3}$. By Lemma~\ref{lem:chernoff},
$2^k-|V_k|\le2^ke^{-c(a)k}$. So the lower density is at least $1-e^{-c(a)k}$ for every $k$, hence equal to $1$, and the density exists and is $1$.
\end{proof}

The bound $e^{-0.0347k}$ is far from the truth for small $k$, as Figure~\ref{fig:terras} shows: only $0.58\%$ of the residues modulo $2^{40}$ have not
stopped after $40$ steps. Notice also what the theorem does not say. A set of density $0$ can be infinite, and the conjecture is about \emph{all}
$n$. Even the statement ``almost every orbit eventually goes below $\frac n{1000}$'' follows the same way, but ``almost every orbit reaches $1$'' does not
follow from any finite number of these steps, because each step loses a set of small but positive density.

\begin{figure}[ht]
\centering
\fig{ch06-terras}
\caption{The proportion $1-\frac{|V_k|}{2^k}$ of residues modulo $2^k$ whose parity vector has not stopped after $k$ steps (no $j\le k$ with
$3^{s_j}<2^j$), on a logarithmic scale, against the Chernoff bound $e^{-0.0347k}$ of Lemma~\ref{lem:chernoff}. The exact values are $0.0625$ at
$k=10$, $0.0058$ at $k=40$ and $0.0019$ at $k=60$.}
\label{fig:terras}
\end{figure}

\begin{remark}[Tao's theorem]\label{rem:tao}
In 2019 Terence Tao proved the strongest result of this kind: for \emph{any} function $f$ with $f(N)\to\infty$, almost every $N$ (in the sense of
logarithmic density) has an orbit that goes below $f(N)$. So almost every orbit gets down to $\log\log\log\log N$, say. The new idea is to show that the
iterates, reduced modulo powers of $3$, stay close to uniformly distributed, by an estimate on the characteristic function of certain random variables,
and to apply this repeatedly as the orbit descends from size $N$ to size $N^{\alpha}$ with $\alpha<1$. The method stops just short of ``almost every orbit reaches $1$''.
Earlier, Ilia Krasikov and Lagarias (2003) had proved that at least $x^{0.84}$ of the integers up to $x$ reach $1$, for large $x$.
\end{remark}

\section{Cycles}

A counterexample would be either an orbit that grows without bound or a cycle other than $1\to2\to1$. Cycles can be attacked by an exact equation.

\begin{proposition}[the cycle equation]\label{prop:cycle}
Let $n_0\to n_1\to\dots\to n_k=n_0$ be a cycle of $T$ of length $k$, with $s$ odd elements among $n_0,\dots,n_{k-1}$. Then $1\le s<k$,
\[
n_0(2^k-3^s)=\rho_k(n_0)\qquad\text{and}\qquad\frac{2^k}{3^s}=\prod_{n_i\text{ odd}}\Bigl(1+\frac1{3n_i}\Bigr).
\]
\end{proposition}
\begin{proof}
If all elements were even, each would be half of the previous one and the cycle could not close; so $s\ge1$. The first equation is
Lemma~\ref{lem:affine} with $T^k(n_0)=n_0$. For the second, multiply the ratios $\frac{n_{i+1}}{n_i}$ around the cycle; the product is $1$. A ratio is
$\frac12$ if $n_i$ is even and $\frac{3n_i+1}{2n_i}=\frac32\bigl(1+\frac1{3n_i}\bigr)$ if $n_i$ is odd, so $1=\frac{3^s}{2^k}\prod_{n_i\text{
odd}}(1+\frac1{3n_i})$. The product is larger than $1$, so $2^k>3^s$ and in particular $k>s$.
\end{proof}

For the trivial cycle, $k=2$, $s=1$ and $\frac43=1+\frac13$. A large cycle must have $\frac{2^k}{3^s}$ just above $1$, that is $\frac ks$ just above
$\log_23$, and the verification limit says how close. To exploit this we need one fact about fractions.

\begin{lemma}[Farey neighbours]\label{lem:farey}
Let $\frac pq<\frac{p'}{q'}$ be fractions with positive denominators and $p'q-pq'=1$. Every fraction $\frac ks$ with $s\ge1$ and
$\frac pq<\frac ks<\frac{p'}{q'}$ has $s\ge q+q'$.
\end{lemma}
\begin{proof}
The numerators $kq-ps$ and $p's-kq'$ are positive integers, so $\frac ks-\frac pq\ge\frac1{sq}$ and $\frac{p'}{q'}-\frac ks\ge\frac1{sq'}$. Adding,
$\frac1{qq'}=\frac{p'}{q'}-\frac pq\ge\frac1{sq}+\frac1{sq'}=\frac{q+q'}{sqq'}$, so $s\ge q+q'$.
\end{proof}

Pairs with $p'q-pq'=1$ are produced by the \emph{Stern--Brocot} descent: start with $\frac11<\alpha<\frac21$, and repeatedly replace one end of the
interval by the mediant $\frac{p+p'}{q+q'}$, keeping $\alpha$ inside. The mediant of a pair with $p'q-pq'=1$ forms such a pair with each end, because
$(p+p')q-p(q+q')=p'q-pq'=1$, and similarly on the other side.

\begin{theorem}[cycles are long]\label{thm:cyclelength}
Suppose every integer $n$ with $1\le n<M$ has an orbit that reaches $1$. Then every cycle of $T$ other than $\{1,2\}$ consists of numbers $\ge M$, and its
length $k$ and number $s$ of odd elements satisfy
\[
0<\frac ks-\log_23<\frac1{3M\log2}.
\]
Consequently, if $\frac pq<\log_23<\frac{p'}{q'}$ with $p'q-pq'=1$ and $\frac{p'}{q'}\ge\log_23+\frac1{3M\log2}$, then $s\ge q+q'$ and
$k>(q+q')\log_23$.
\end{theorem}
\begin{proof}
An element of another cycle cannot reach $1$, since its orbit stays in the cycle; so it is at least $M$. By Proposition~\ref{prop:cycle},
$1<\frac{2^k}{3^s}\le\bigl(1+\frac1{3M}\bigr)^s$. Taking logarithms and using $\log(1+u)<u$,
$0<k\log2-s\log3<\frac s{3M}$; divide by $s\log2$. For the second part, $\frac ks$ lies strictly between $\log_23$ and
$\log_23+\frac1{3M\log2}\le\frac{p'}{q'}$, hence strictly between $\frac pq$ and $\frac{p'}{q'}$; apply Lemma~\ref{lem:farey}.
\end{proof}

\begin{corollary}\label{cor:cyclebound}
Every cycle of $T$ other than $\{1,2\}$ has at least $72\,057\,431\,991$ odd elements and length at least $114\,208\,327\,604$. In terms of the original map $C$,
whose cycle has $k+s$ steps, its length is at least $186\,265\,759\,595$.
\end{corollary}
\begin{proof}
Take $M=2^{71}$ (Barina's verification), so $\frac1{3M\log2}=2.04\cdot10^{-22}$. The Stern--Brocot descent towards $\log_23=1.58496\ldots$, computed
with exact fractions and $80$ digits for $\log_23$ (Python companion, P6.5), stops with
\[
\frac pq=\frac{103\,768\,467\,013}{65\,470\,613\,321},\qquad\frac{p'}{q'}=\frac{10\,439\,860\,591}{6\,586\,818\,670},
\]
where $p'q-pq'=1$, $\frac pq<\log_23$, and $\frac{p'}{q'}-\log_23=2.2\cdot10^{-21}$, larger than $2.04\cdot10^{-22}$. Their mediant
$\frac{114\,208\,327\,604}{72\,057\,431\,991}$ lies within $1.1\cdot10^{-22}$ above $\log_23$, which is why the descent cannot go further. So
$s\ge q+q'=72\,057\,431\,991$, and $k>s\log_23$ gives $k\ge114\,208\,327\,604$.
\end{proof}

The same computation with $M=2^{68}$ gives the same numbers, and with $M=2^{40}$, the limit known in 1993, it gives $k\ge17\,087\,915$, the bound found
then by Shalom Eliahou. Sharper versions of the argument, which use more of the cycle equation than its smallest element, give larger bounds.

\begin{remark}[cycles of a given shape]\label{rem:mcycles}
A cycle made of one run of odd steps followed by one run of halvings is called a $1$-cycle (or circuit). Ray Steiner proved in 1977 that the only
$1$-cycle is the trivial one, using Alan Baker's lower bounds for linear forms in logarithms, here for $|k\log2-s\log3|$; John Simons and Benne de
Weger (2005) excluded cycles made of up to $68$ such blocks, and Christian Hercher (2023) up to $91$. Baker's theorem is what replaces our Farey lemma
when nothing is known about the size of the elements.
\end{remark}

\section{Why it is hard}

\paragraph{The neighbours have cycles.} The map $n\mapsto3n-1$ (for odd $n$, with halving for even $n$) has the same coin model and the same negative
drift, and it has three known cycles, through $1$, $5$ and $17$: for instance $5\to14\to7\to20\to10\to5$ (Exercise~\ref{ex:minusone}). It is the
$3n+1$ map on the negative integers in disguise. The map $5n+1$ has cycles through $13$ and $17$, and its drift is positive, so most of its orbits are
expected to escape to infinity (Exercise~\ref{ex:fiveone}). Any proof of the Collatz conjecture must therefore use something that distinguishes $3n+1$
on the positive integers from these close relatives, and the coin model, which is the same for all of them, cannot.

\paragraph{Base 2 against base 3.} The halvings act on the binary digits of $n$, the multiplications by $3$ on its ternary digits, and the walk mixes
the two. Theorem~\ref{thm:paritybij} is the precise sense in which the mixing is perfect for the first $k$ steps, and it is exactly this perfection that
leaves nothing to hold on to: there is no quantity that decreases along every orbit.

\paragraph{Conway's theorem.} In 1972 John Conway considered all maps of the form $g(n)=a_in+b_i$ for $n\equiv i\pmod P$, with rational $a_i,b_i$ such that
$g$ maps integers to integers; the Collatz map is the case $P=2$. He proved that such maps can simulate any computer program, so that no algorithm can
decide, for every map of this kind and every starting value, whether the orbit ever reaches a power of $2$. A proof of the Collatz conjecture must use
something special about this one map. In 2007 Stuart Kurtz and Janos Simon showed that the general question ``does every orbit of $g$ reach $1$?'' is
even harder than the halting problem.

\begin{remark}[FRACTRAN]\label{rem:fractran}
Conway later turned this into a programming language. A FRACTRAN program is a list of fractions $f_1,\dots,f_r$; from the current integer $n$, the next
one is $nf_i$ for the first $i$ such that $nf_i$ is an integer, and the program stops if there is none. The exponents of $n=2^a3^b5^c\cdots$ are the
registers of a machine, and a fraction tests and changes them. Conway's PRIMEGAME is the list
\[
\tfrac{17}{91},\tfrac{78}{85},\tfrac{19}{51},\tfrac{23}{38},\tfrac{29}{33},\tfrac{77}{29},\tfrac{95}{23},\tfrac{77}{19},\tfrac1{17},\tfrac{11}{13},\tfrac{13}{11},\tfrac{15}2,\tfrac17,\tfrac{55}1;
\]
started at $2$, the powers of $2$ that appear are exactly $2^p$ for the primes $p$, in increasing order. The Python companion runs it (P6.6). Each
FRACTRAN step is a map of Conway's type, with $P$ the least common multiple of the denominators.
\end{remark}

\section{Reading a paper: Lagarias and Tao}

Jeffrey Lagarias, \emph{The $3x+1$ problem and its generalizations} (American Mathematical Monthly, 1985), is the survey from which every later
paper starts. Read the introduction for the history (our account of the names follows it) and the part on stopping times for Terras's theorem, which it
proves with the parity vectors of Theorem~\ref{thm:paritybij}; skim the generalisations and the list of open problems. The idea worth stealing is the coding of an orbit by
its parity vector, which turns a dynamical question into a counting question modulo $2^k$.

Terence Tao, \emph{Almost all orbits of the Collatz map attain almost bounded values} (Forum of Mathematics, Pi, 2022; arXiv 1909.03562), is long, but
its introduction is a model. Read the introduction for the statement, the history and the heuristic: the orbit is followed through the
\emph{Syracuse map} on odd numbers, the exponents of $2$ removed at each step are approximately independent geometric random variables, and the heart
of the proof is that the iterates, reduced modulo $3^n$, become close to uniformly distributed. The idea worth stealing is to prove a statement for
almost all starting points by propagating a distribution, not a single orbit.

\section{Open doors}

\begin{itemize}
\item Compute the stopping-time records of $T$ up to $10^9$. Plot $\sigma(n)$ for the records against $\log n$. Does a straight line fit?
\item For odd $d$ not divisible by $3$, the map $n\mapsto\frac{3n+d}2$ (odd $n$), $\frac n2$ (even $n$) has cycles that depend on $d$. List them for
$d\le100$ and starting values up to $10^6$. Which $d$ give the most cycles?
\item Use Lemma~\ref{lem:affine} to find all residues modulo $2^{20}$ that do not stop within $20$ steps, and check that their number is $27\,328$, as in
Figure~\ref{fig:terras}.
\item Improve Theorem~\ref{thm:cyclelength}: use that the smallest element of a cycle is odd and at least $M$, and that every odd element is followed by a
larger one, to sharpen the inequality on $\frac{2^k}{3^s}$. How much does the bound on $s$ improve?
\item For the $5n+1$ map, estimate the proportion of $n\le10^5$ whose orbit exceeds $10^{30}$ without entering a known cycle.
\end{itemize}

\section{Exercises}

\begin{exercise}[$\star$]\label{ex:stopsmall}
Show that $\sigma(n)=1$ if and only if $n$ is even, that $\sigma(n)=2$ if and only if $n\equiv1\pmod4$ and $n>1$, that no $n$ has $\sigma(n)=3$, and that
$\sigma(n)=4$ if and only if $n\equiv3\pmod{16}$. What are the densities of these four sets?
\end{exercise}

\begin{exercise}[$\star$]\label{ex:climb}
Show that $T(n)+1\le\frac32(n+1)$, with equality for odd $n$. Deduce that $T^j(2^k-1)=3^j2^{k-j}-1$ for $0\le j\le k$, and that $\sigma(2^k-1)>k$. So there
are numbers with arbitrarily long climbs.
\end{exercise}

\begin{exercise}[$\star$]\label{ex:backward}
Show that $m$ has a predecessor under $C$ other than $2m$ if and only if $m\equiv4\pmod6$, and that this predecessor is then the odd number $\frac{m-1}3$.
Show that if $3\mid m$, the part of the backward tree above $m$ is the single chain $m,2m,4m,\dots$
\end{exercise}

\begin{exercise}[$\star\star$]\label{ex:oneodd}
Using Lemma~\ref{lem:affine}, show that the only cycle of $T$ with exactly one odd element is $\{1,2\}$.
\end{exercise}

\begin{exercise}[$\star\star$]\label{ex:cyclesmall}
Assume that every $n<100$ reaches $1$ (check a few by hand). Run the Stern--Brocot descent of Theorem~\ref{thm:cyclelength} by hand with $M=100$, and deduce
that a non-trivial cycle of $T$ has $s\ge17$ and $k\ge27$.
\end{exercise}

\begin{exercise}[$\star$]\label{ex:minusone}
Check the cycles $5\to14\to7\to20\to10\to5$ and the cycle through $17$ for the map $n\mapsto3n-1$ (odd $n$), $n\mapsto\frac n2$ (even $n$). Show that
$n\mapsto-n$ conjugates this map on the positive integers to the Collatz map on the negative integers. What does this say about a possible proof?
\end{exercise}

\begin{exercise}[$\star\star$]\label{ex:fiveone}
For the map $n\mapsto5n+1$ (odd $n$), $n\mapsto\frac n2$ (even $n$), find the cycles through $13$ and $17$. Show that Lemmas~\ref{lem:affine}
and~\ref{lem:shift} and Theorem~\ref{thm:paritybij} hold for its shortcut map $\frac{5n+1}2$, $\frac n2$, with $3$ replaced by $5$. Compute the drift of the
coin model. Which step of the proof of Theorem~\ref{thm:terras} fails?
\end{exercise}

\begin{exercise}[$\star$]\label{ex:drift}
Justify the constants $6.95$ and $10.43$ of the heuristic, and show that $10.43\ldots=\frac3{\log(4/3)}$. Explain this second form directly from the map
$C$: what does one odd step and the halvings that follow it do to $n$ on average, and how many steps of $C$ does that take?
\end{exercise}

\begin{exercise}[$\star$]\label{ex:fractran}
Show that the FRACTRAN program $\frac32$ turns $2^a3^b$ into $3^{a+b}$, and that the program $\frac5{6}$ turns $2^a3^b$ into
$2^{a-m}3^{b-m}5^m$ with $m=\min(a,b)$. So FRACTRAN can add and take minima.
\end{exercise}

\begin{exercise}[$\star\star$]\label{ex:chernoffsharp}
Let $0<p<1$ and $j=pk$ an integer. Show that $\binom kjp^j(1-p)^{k-j}\ge\frac1{k+1}$, by comparing it with the other terms of $(p+(1-p))^k$, given that it is
the largest. Deduce that $\frac1k\log\bigl(2^{-k}\sum_{s\ge ak}\binom ks\bigr)\to-c(a)$ for $\frac12<a<1$, so the exponent of Lemma~\ref{lem:chernoff} is
the right one.
\end{exercise}

\section{Hints}

\begin{hint}{ex:stopsmall}
Write $n=4m+1$, then $n=16m+r$ for $r=3,7,11,15$, and compute the first four iterates.
\end{hint}
\begin{hint}{ex:climb}
For odd $n$, $T(n)+1=\frac{3(n+1)}2$. Apply this $j$ times to $n=2^k-1$, as long as the iterate is odd.
\end{hint}
\begin{hint}{ex:backward}
A predecessor $p$ of $m$ is either $2m$ or an odd $p$ with $3p+1=m$. Powers of $2$ times a multiple of $3$ are multiples of $3$.
\end{hint}
\begin{hint}{ex:oneodd}
With one odd step at position $j$, $\rho_k=2^j$. The cycle equation becomes $n(2^k-3)=2^j$.
\end{hint}
\begin{hint}{ex:cyclesmall}
$\frac1{300\log2}=0.00481$. Start from $\frac11,\frac21$ and look at $\frac32,\frac53,\frac85,\dots$
\end{hint}
\begin{hint}{ex:minusone}
$3(-n)+1=-(3n-1)$.
\end{hint}
\begin{hint}{ex:fiveone}
The proofs of the lemmas only use that $3$ is odd. For the theorem, a vector stops when $5^{s_j}<2^j$.
\end{hint}
\begin{hint}{ex:drift}
Mean step of $\log T^k(n)$ per step of $T$, and mean cost of a step of $T$ in steps of $C$.
\end{hint}
\begin{hint}{ex:fractran}
In both programs only one fraction exists, so it is applied as long as the result is an integer.
\end{hint}
\begin{hint}{ex:chernoffsharp}
The binomial expansion of $(p+(1-p))^k=1$ has $k+1$ terms. And $p^j(1-p)^{k-j}=e^{-k(-p\log p-(1-p)\log(1-p))}$.
\end{hint}

\section{Solutions}

\begin{solution}{ex:stopsmall}
\idea Each congruence class modulo $2^k$ fixes the first $k$ steps (Lemma~\ref{lem:shift}); compute them.
\writeup An even $n$ has $T(n)=\frac n2<n$; an odd $n$ has $T(n)=\frac{3n+1}2>n$. So $\sigma(n)=1$ exactly for even $n$. For $n=4m+1$ with $m\ge1$:
$T(n)=6m+2$, $T^2(n)=3m+1<n$, so $\sigma(n)=2$; for $n=1$, $\sigma=\infty$. For $n=4m+3$: $T(n)=6m+5$ is odd, so $T^2(n)=9m+8>n$, and $T^3(n)\ge\frac{9m+8}2>n$.
So $\sigma(n)\ge4$ for $n\equiv3\pmod4$, and $\sigma(n)=3$ never happens. Modulo $16$:
\begin{itemize}
\item $16m+3\to24m+5\to36m+8\to18m+4\to9m+2<n$, so $\sigma=4$;
\item $16m+7\to24m+11\to36m+17\to54m+26\to27m+13>n$;
\item $16m+11\to24m+17\to36m+26\to18m+13\to27m+20>n$;
\item $16m+15\to24m+23\to36m+35\to54m+53\to81m+80>n$.
\end{itemize}
So $\sigma(n)=4$ exactly for $n\equiv3\pmod{16}$. The densities are $\frac12$, $\frac14$, $0$ and $\frac1{16}$.
\end{solution}

\begin{solution}{ex:climb}
\idea $n+1$ is multiplied by exactly $\frac32$ at an odd step.
\writeup For even $n$, $T(n)+1=\frac n2+1\le\frac32(n+1)$. For odd $n$, $T(n)+1=\frac{3n+3}2=\frac32(n+1)$. Now let $n=2^k-1$ and prove
$T^j(n)=3^j2^{k-j}-1$ by induction on $j\le k$: it holds for $j=0$, and if $j<k$ then $3^j2^{k-j}-1$ is odd, so
$T^{j+1}(n)+1=\frac32\cdot3^j2^{k-j}=3^{j+1}2^{k-j-1}$. For $1\le j\le k$, $T^j(n)=3^j2^{k-j}-1>2^k-1$, so $\sigma(n)>k$.
\end{solution}

\begin{solution}{ex:backward}
\idea Solve $C(p)=m$.
\writeup $C(p)=m$ means $p=2m$ (if $p$ is even) or $3p+1=m$ with $p$ odd. The second needs $m\equiv1\pmod3$ and $p=\frac{m-1}3$ odd, that is
$m-1\equiv3\pmod6$, or $m\equiv4\pmod6$; it also needs $p\ge1$, true for $m\ge4$. If $3\mid m$, every $2^jm$ is a multiple of $3$, never
$\equiv4\pmod6$, so each $2^jm$ has the single predecessor $2^{j+1}m$, and the tree above $m$ is a chain.
\end{solution}

\begin{solution}{ex:oneodd}
\idea Lemma~\ref{lem:affine} computes $\rho_k$ exactly when there is one odd step.
\writeup Let $n$ be an element of the cycle, of length $k$, and let $j$ be the position of the odd step, $0\le j\le k-1$. By the recursion of
Lemma~\ref{lem:affine}, $\rho$ stays $0$ until step $j$, becomes $3\cdot0+2^j=2^j$, and stays $2^j$ through the halvings. So $T^k(n)=n$ reads
$2^kn=3n+2^j$, that is $n(2^k-3)=2^j$. The odd number $2^k-3$ divides $2^j$, so $2^k-3=1$, $k=2$, and $n=2^j\in\{1,2\}$. The cycle is $\{1,2\}$.
\end{solution}

\begin{solution}{ex:cyclesmall}
\idea Run the descent, keeping the right end at least $\log_23+\delta$.
\writeup Here $\delta=\frac1{300\log2}=0.00481$ and $\log_23=1.58496$, so the right end must stay $\ge1.58977$. From $\frac11<\log_23<\frac21$: the mediant
$\frac32=1.5$ is below, so it becomes the left end; $\frac53=1.667$ and then $\frac85=1.6$ are above $1.58977$, so they become the right end;
$\frac{11}7=1.5714$ and $\frac{19}{12}=1.5833$ are below, left end. The next mediant $\frac{27}{17}=1.58824$ is above $\log_23$ but below $1.58977$, so the
descent stops with $\frac pq=\frac{19}{12}$, $\frac{p'}{q'}=\frac85$, and indeed $8\cdot12-19\cdot5=1$. By Theorem~\ref{thm:cyclelength}, $s\ge12+5=17$ and
$k>17\log_23=26.94$, so $k\ge27$. (Checking that all $n<100$ reach $1$ is a matter of minutes by hand with the stopping-time trick of
Proposition~\ref{prop:stopping}, or seconds by computer.)
\end{solution}

\begin{solution}{ex:minusone}
\idea Negate.
\writeup $5\to14\to7\to20\to10\to5$, and $17\to50\to25\to74\to37\to110\to55\to164\to82\to41\to122\to61\to182\to91\to272\to136\to68\to34\to17$. For the
Collatz map at $-n$: if $n$ is odd, $3(-n)+1=-(3n-1)$, and if $n$ is even, $\frac{-n}2=-\frac n2$. So $C(-n)=-g(n)$ where $g$ is the $3n-1$ map, and
the cycles above are the cycles $-5\to-14\to\dots$ and $-17\to\dots$ of the Collatz map on negative integers. A proof of the conjecture must use the
positivity of $n$ somewhere; an argument that works the same way for all integers is doomed.
\end{solution}

\begin{solution}{ex:fiveone}
\idea The coin model is the same, but the drift has the wrong sign.
\writeup $13\to66\to33\to166\to83\to416\to208\to104\to52\to26\to13$ and $17\to86\to43\to216\to108\to54\to27\to136\to68\to34\to17$. In the proofs of
Lemmas~\ref{lem:affine} and~\ref{lem:shift}, replace $3$ by $5$: the only property used is that $5^{s}2^{k-i}$ is even when $k>i$ and $5^s$ is odd, so
the proofs go through, and so does Theorem~\ref{thm:paritybij}. The drift is $\frac12\log\frac52+\frac12\log\frac12=\frac12\log\frac54>0$. In the proof
of Theorem~\ref{thm:terras}, a vector stops when $5^{s_j}<2^j$, that is $s_j<0.4307j$, fewer than half the steps odd: the proportion of vectors that
stop by time $k$ no longer tends to $1$, and the argument gives only a lower density bounded away from $1$. (In fact one expects most orbits of $5n+1$ to
grow without bound, but no orbit is known to do so.)
\end{solution}

\begin{solution}{ex:drift}
\idea Divide the distance to travel, $\log n$, by the mean step.
\writeup By Corollary~\ref{cor:coins} the mean step of $\log T^k(n)$ is $\frac12\log\frac32+\frac12\log\frac12=\log\frac{\sqrt3}2=-0.14384$, so reaching $1$
takes about $\frac{\log n}{0.14384}=6.952\log n$ steps of $T$. A step of $T$ is one step of $C$ if it halves and two if it is odd, $1.5$ on average, so
$10.428\log n$ steps of $C$. Directly: after an odd step $3n+1\approx3n$, the number of halvings that follow is $j\ge1$ with probability $2^{-j}$, mean
$2$; so one ``odd block'' multiplies $n$ by $\frac34$ on average in the logarithmic sense ($\log3-2\log2=\log\frac34$) and costs $3$ steps of $C$.
Reaching $1$ takes $\frac{\log n}{\log(4/3)}$ blocks, so $\frac{3\log n}{\log(4/3)}=10.428\log n$ steps.
\end{solution}

\begin{solution}{ex:fractran}
\idea A single fraction is applied while it can be.
\writeup With $\frac32$: $2^a3^b\cdot\frac32=2^{a-1}3^{b+1}$ as long as $a\ge1$; after $a$ steps we reach $3^{a+b}$, and $3^{a+b}\cdot\frac32$ is not an
integer, so the program stops. With $\frac56$: $2^a3^b\cdot\frac56=2^{a-1}3^{b-1}5$ as long as $a,b\ge1$; after $m=\min(a,b)$ steps one of the
exponents is $0$, the product is no longer an integer, and the program stops at $2^{a-m}3^{b-m}5^m$.
\end{solution}

\begin{solution}{ex:chernoffsharp}
\idea The largest of $k+1$ nonnegative terms with sum $1$ is at least $\frac1{k+1}$.
\writeup The terms $t_i=\binom kip^i(1-p)^{k-i}$, $0\le i\le k$, sum to $1$. The ratio $\frac{t_{i+1}}{t_i}=\frac{(k-i)p}{(i+1)(1-p)}$ is at least $1$
exactly when $i+1\le(k+1)p$, so $t_i$ increases up to $i\approx kp$ and then decreases; with $j=pk$ an integer, $t_j$ is the largest term (compare
$t_{j}$ with $t_{j\pm1}$ using the ratio). So $t_j\ge\frac1{k+1}$, that is $\binom kj\ge\frac1{k+1}e^{kH(p)}$ with $H(p)=-p\log p-(1-p)\log(1-p)$. Now let
$\frac12<a<1$ and $j=\lceil ak\rceil$, $p=\frac jk\to a$. Then $\sum_{s\ge ak}\binom ks\ge\binom kj\ge\frac1{k+1}e^{kH(j/k)}$, so
$\frac1k\log\bigl(2^{-k}\sum_{s\ge ak}\binom ks\bigr)\ge H(\tfrac jk)-\log2-\frac{\log(k+1)}k\to H(a)-\log2=-c(a)$, by continuity of $H$. With
Lemma~\ref{lem:chernoff} the limit is exactly $-c(a)$.
\end{solution}

\section*{Sources and further reading}

\begin{itemize}
\item J.~C.~Lagarias, \emph{The $3x+1$ problem and its generalizations}, American Mathematical Monthly 92 (1985), 3--23 (the history of the problem,
Kakutani's recollection quoted in the epigraph, and Terras's theorem). J.~C.~Lagarias (ed.), \emph{The Ultimate Challenge: The $3x+1$ Problem}, American
Mathematical Society, 2010.
\item R.~K.~Guy, \emph{Unsolved Problems in Number Theory}, 3rd ed., Springer, 2004, Problem E16 (with Erd\H{o}s's remark).
\item R.~Terras, \emph{A stopping time problem on the positive integers}, Acta Arithmetica 30 (1976), 241--252. C.~J.~Everett, \emph{Iteration of the
number-theoretic function $f(2n)=n$, $f(2n+1)=3n+2$}, Advances in Mathematics 25 (1977), 42--45.
\item D.~Barina, \emph{Convergence verification of the Collatz problem}, The Journal of Supercomputing 77 (2021), 2681--2688 (all $n<2^{68}$).
D.~Barina, \emph{Improved verification limit for the convergence of the Collatz conjecture}, The Journal of Supercomputing 81 (2025), article 810
(all $n<2^{71}$).
\item S.~Eliahou, \emph{The $3x+1$ problem: new lower bounds on nontrivial cycle lengths}, Discrete Mathematics 118 (1993), 45--56.
\item R.~P.~Steiner, \emph{A theorem on the Syracuse problem}, Proceedings of the 7th Manitoba Conference on Numerical Mathematics (1977), 553--559.
J.~Simons and B.~de Weger, \emph{Theoretical and computational bounds for $m$-cycles of the $3n+1$ problem}, Acta Arithmetica 117 (2005), 51--70.
C.~Hercher, \emph{There are no Collatz $m$-cycles with $m\le91$}, Journal of Integer Sequences 26 (2023), Article 23.3.5.
\item I.~Krasikov and J.~C.~Lagarias, \emph{Bounds for the $3x+1$ problem using difference inequalities}, Acta Arithmetica 109 (2003), 237--258.
\item T.~Tao, \emph{Almost all orbits of the Collatz map attain almost bounded values}, Forum of Mathematics, Pi 10 (2022), e12; arXiv:1909.03562.
\item J.~H.~Conway, \emph{Unpredictable iterations}, Proceedings of the 1972 Number Theory Conference, University of Colorado, Boulder (1972), 49--52.
J.~H.~Conway, \emph{FRACTRAN: a simple universal programming language for arithmetic}, in T.~M.~Cover and B.~Gopinath (eds.), \emph{Open Problems in
Communication and Computation}, Springer, 1987, 4--26. S.~A.~Kurtz and J.~Simon, \emph{The undecidability of the generalized Collatz problem}, in
\emph{Theory and Applications of Models of Computation} (TAMC 2007), Lecture Notes in Computer Science 4484, Springer, 2007, 542--553.
\end{itemize}
